\begin{figure}[htbp]
\centering
\begin{tikzpicture}[
  node distance=6mm,
  caja/.style={draw, rounded corners=2pt, align=center, font=\small, text width=0.50\linewidth, minimum height=9mm, inner sep=4pt},
  clave/.style={caja, line width=1pt, fill=gray!15},
  lado/.style={draw, dashed, rounded corners=2pt, align=center, font=\footnotesize, text width=0.28\linewidth, inner sep=4pt},
  flecha/.style={-{Stealth[length=2.4mm]}, thick}]
  \node[caja] (tc) {TC receptora sin osteosíntesis\\ (subconjunto local de CTPelvic1K)};
  \node[caja, below=of tc] (mask) {Máscaras anatómicas (TotalSegmentator) y medición\\ del corredor óseo (marco de Kaiser et al.)};
  \node[clave, below=of mask] (muest) {\textbf{Muestreador de colocación} (Objetivo~2)\\ poses perturbadas alrededor del eje del corredor;\\ SAP frente a las dos distribuciones de la referencia clínica};
  \node[caja, below=of muest] (tor) {Tornillo paramétrico en la pose muestreada:\\ máscara $M$ y región $G = M \cup B_{\delta}$};
  \node[clave, below=of tor] (sint) {\textbf{Sintetizador} (Objetivo~3)\\ \emph{inpainting} por difusión en el dominio de imagen dentro de $G$;\\ fuera de $G$ se copia el volumen de origen};
  \node[caja, below=of sint] (eva) {Evaluación de apariencia: métricas de Peters et al.\\ frente a copia y pegado y al protocolo físico;\\ reconstrucción de implantes reales};
  \node[lado, right=10mm of tor] (comp) {\emph{Compuerta de representación} (Objetivo~1), previa a la cadena: su resultado negativo descartó la ruta latente y fijó el dominio de imagen};
  \draw[flecha] (tc) -- (mask);
  \draw[flecha] (mask) -- (muest);
  \draw[flecha] (muest) -- (tor);
  \draw[flecha] (tor) -- (sint);
  \draw[flecha] (sint) -- (eva);
  \draw[flecha, dashed] (comp.south) |- ([yshift=1mm]sint.east);
\end{tikzpicture}
\caption{Cadena propuesta. Las cajas con borde grueso son los dos componentes que se diseñan y se evalúan por separado; la caja punteada es la compuerta que fijó el dominio del sintetizador.}
\label{fig:pipeline}
\end{figure}
